\scanpage{36}
N.B. Fixons un pt \(y_0\) de \(Y\). Alors le gr des familles algébriques \(Y(k)\mathbin{-}A^i_{\mathrm{alb}}(X)\) est isom à \(A^i_{\mathrm{alb}}(X)\times\operatorname{Hom\,alg}(A_Y,A^i_{\mathrm{alb}}(X))\), où le dernier facteur désigne les familles algébriques qui sont des hom de groupes, i.e. nulles à l'origine. On peut donc se limiter à étudier les hom ``alg.'' \(A_{\deleted{\unsure{Y}}}(k)\mathbin{-}\figure{\begin{tikzpicture}[baseline=(target.base)]\node[draw,ellipse,inner sep=2pt] (target) {$A^i_{\mathrm{alb}}(X)$};\end{tikzpicture}}\), \inserted{où \(A\) est une VA,} et, interprétant \(A\) comme
\[
A_{0\,\mathrm{alb}}(A)=A^\alpha_{\mathrm{alb}}(A)^{(\mathrm{alg.\,équiv.\,0})}\qquad(\alpha=\dim A)
\]
on peut considérer \deleted{\unsure{le}} ces objets comme des cas particuliers des \uline{homs algébriques}
\[
\varphi:P^j(X')\mathbin{-}P^i(X).
\]
\margin{\diagram{\begin{tikzcd}[ampersand replacement=\&]A(k)\simeq P^{\unsure{\alpha|n}}(A)\arrow[r]\&P^i(X)\end{tikzcd}}}
Un hom de groupes \(\varphi\) est dit ``algébrique'' s'il existe un entier \(N>0\) tel que \(N\varphi\) soit définissable pr un cycle algébrique sur \(\unsure{X\times Y'|X\times X'}\).

(Alors \(\varphi\) sera définissable pr un cycle alg. à coeff \uline{rationnels} --- mais attention, un cycle à coeff rationnels \uline{quelconque} ne définit pas en général un hom \(P^j(X')\mathbin{-}P^i(X)\), \(=\) si les degrés sont les bons\ldots.). Ces homs \uline{se composent}
\mspage{35}
